\section{Construction of $Q_{n,k}$}\label{section: coupling}

\subsection{Odd $k$-cycle}
Let $t^* = \lfloor  \frac{n\log n}{k}\rfloor - \frac{n}{2k}\log\log \log n+\frac{4n}{k}\log\log\log \log n$. Let $S_{t^*}$ denote the set of elements in $[n]$ that are not touched by time $t^*$. We consider the following two marking schemes. \JC{Motivation?}

\noindent\textbf{Marking scheme A:} at time $t^*$, mark all the elements in $[n] \setminus S_{t^*}$. Let $S_t$ denote the set of all the unmarked elements at time $t$ and let $M_t = [n] \setminus S_t$. At time $t + 1$, we choose a random k-cycle $\sigma = (i_1 \cdots i_k)$ according to the distribution $Q_{n,k}$ defined as \JC{mention that in this case $i_1,...,i_k$ can take different values while in $P_{n,k}$ they have to be distinct}
    \[
Q_{n,k}=\begin{cases}
(i\cdots i) & \text{with probability } C_tq_t \ \text{for } i\in S_{t}\\[6pt]
(i_1 \cdots i_k) & \text{with probability } q_t\ \text{for } i_1 \in S_{t}, i_2, \cdots, i_k \in [n] \setminus S_{t}\\[6pt]
(i_1 \cdots i_k) & \text{with probability } \left( \frac{n!}{k(n-k)!} \right)^{-1} \ \text{otherwise},
\end{cases}
\]
where $C_t = (|M_t|-1)_{k-2}$ and $q_t = \frac{|M_t|}{|M_t|+1}\left(\frac{n!}{k(n-k)!}\right)^{-1}$.

We update the marked set according to the following rule.
\begin{itemize}
    \item If $\sigma=(i\cdots i)$ with $i\in S_t$, then we mark $i$.
    \item If $\sigma=(i_1\cdots i_k)$ with $i_1\in S_t$ and
$i_2,\ldots,i_k\in M_t$, then we mark $i_1$.
\item In all remaining cases, namely when the cycle contains more than one unmarked element, we update the marked set by pulling it back under $\sigma$:
\[
    M_{t+1}
    =
    \{i\in[n]:\sigma(i)\in M_t\}
    =
    \sigma^{-1}(M_t),
\]
and hence
\[
    S_{t+1}
    =
    \{i\in[n]:\sigma(i)\in S_t\}
    =
    \sigma^{-1}(S_t).
\]
\end{itemize}




\textcolor{red}{Almost same as scheme A}
\noindent\textbf{Marking scheme B:} at time $t^*$, mark all the elements in $[n] \setminus S_{t^*}$. Let $S_t$ denote the set of all the unmarked elements at time $t$ and let $M_t = [n] \setminus S_t$. At time $t + 1$, we choose a random $k$-cycle $\sigma = (i_1 \cdots i_k)$ according to the distribution $P_{n,k}$. We update the marked set according to the following rule.
\begin{itemize}
    \item If $\sigma=(i_1\cdots i_k)$ with $i_1\in S_t$ and
$i_2,\ldots,i_k\in M_t$, then we mark $i_1$.
\item In all remaining cases, namely when the cycle contains more than one unmarked element, we update the marked set by pulling it back under $\sigma$:
\[
    M_{t+1}
    =
    \{i\in[n]:\sigma(i)\in M_t\}
    =
    \sigma^{-1}(M_t),
\]
and hence
\[
    S_{t+1}
    =
    \{i\in[n]:\sigma(i)\in S_t\}
    =
    \sigma^{-1}(S_t).
\]
\end{itemize}

\noindent\textbf{Coupled Markov processes.}
Let
\[
    \rho_{t^*}=Y_{t^*}=Z_{t^*}
\]
be a common random variable sampled uniformly from \(\mathfrak A_{S_{t^*}}\). We define four Markov processes
\[
    (X_t)_{t\geq t^*},\qquad
    (Y_t)_{t\geq t^*},\qquad
    (Z_t)_{t\geq t^*},\qquad
    (\rho_t)_{t\geq t^*}
\]
on a common probability space as follows:

\begin{itemize}
    \item The process $(X_t)_{t\geq t^*}$ is the original Markov chain, evolving according to the transition kernel \(P_n\). Let \(\tau\) denote the first time at which all elements have been touched.
    \item The process $(Z_t)_{t\geq t^*}$ is initialized by $Z_{t^*}$, evolving according to the same transition kernel \(P_n\) as $X_t$. Let \(\tau_1\) denote the first time at which all elements have been marked.
    \item The process $(Y_t)_{t\geq t^*}$ is initialized by $Y_{t^*}$, evolving according to the time-inhomogeneous transition kernel \(Q_{n,t}\). Let \(\tau_2\) denote the first time at which all elements have been marked.
    \item Finally, the process $(\rho_t)_{t\geq t^*}$ is initialized by $\rho_{t^*}$, evolving according to the same \(Q_{n,t}\)-updates as $Y_t$, but with a modified action on the current state. More precisely, at time \(t+1\), sample \(\sigma\sim Q_{n,t}\). If either
\[
    \sigma=(i\cdots i)\quad \text{for some } i\in S_t,
\]
or
\[
    \sigma=(i_1\cdots i_k),\qquad i_1\in S_t,\qquad
    i_2,\ldots,i_k\in [n]\setminus S_t,
\]
then set
\[
    \rho_{t+1}:=\sigma\rho_t.
\]
In all other cases, set
\[
    \rho_{t+1}:=\sigma^{-1}\rho_t\sigma.
\]
Let \(\tau_3\) denote the first time at which all elements have been marked.
\end{itemize}

\subsection{Proof of \Cref{thm: hitting time mixing_main thm}}


We split the proof of \Cref{thm: hitting time mixing_main thm} into two main parts. 


First, let $\mathfrak{A}(M):=\{\sigma\in \mathfrak{A}_n:\sigma(x)=x,\;\forall x\notin M\}$. We show that, conditioned on the marked set \(M_t\), the random permutation \(\rho_t\) is uniform on \(\mathfrak A(M_t)\). In particular, at the marking time \(\tau_2\), all elements are marked, and hence \(\rho_{\tau_2}\) is uniform on \(\mathfrak A_n\).

Second, we compare the auxiliary chains with the original chain. More precisely, we show that
\[
    d_{TV}(\rho_{\tau_3},Y_{\tau_2}),\qquad
    d_{TV}(Y_{\tau_2},Z_{\tau_1}),\qquad
    d_{TV}(Z_{\tau_1},X_{\tau})
\]
are all small. Combining these estimates with the uniformity of \(\rho_{\tau_2}\) yields the desired hitting-time mixing statement.

\begin{lemma}
For $t \geq t^*$, conditioned on the marked set $M_t$, $\rho_t$ is uniform on $\mathfrak{A}(M_t)$. In particular, $\rho_{\tau_3}$ is uniform on $\mathfrak{A}_n$.
\end{lemma}
\begin{proof}
    This is true at time $t^*$ by construction. Suppose it is true at time $t\geq t^*$. At time $t+1$, there are two possibilities. The first is that $|M_{t+1}| = |M_t|$. This corresponds to the third circumstance in our marking scheme, in this case, $\rho_{t+1}=\sigma^{-1}\rho_t\sigma$, and 
        $\operatorname{sgn}(\sigma^{-1}\rho_t\sigma) =\operatorname{sgn}(\rho_t)$. Because $\rho_t$ is uniform on $\mathfrak{A}(M_t) = \{\delta \in \mathfrak{A}_n: \delta(x) = x, \forall x \notin M_t\}$, $\sigma^{-1}\rho_t \sigma$ is uniform on $A = \{\sigma^{-1}\delta\sigma: \delta \in \mathfrak{A}_n, \delta(x) = x, \forall x \notin M_t\}$. For any $\sigma^{-1}\delta\sigma \in A$, since $\sigma^{-1} \delta \sigma \in \mathfrak{A}_n$ and for every $\sigma(x) \notin M_t$, $\sigma^{-1}\delta \sigma (x) = \sigma^{-1}\sigma(x) = x$, we know that $\sigma^{-1}\delta\sigma \in \mathfrak{A}(\sigma^{-1}(M_t))$. On the other hand, for any $\kappa \in \mathfrak{A}(\sigma^{-1}(M_t))$, there exists some $\delta \in \mathfrak{A}_n$ such that $\kappa = \sigma^{-1}\delta\sigma$. We know that $\delta(\sigma(x)) = \sigma(x)$ for all $\sigma(x) \not \in M_t$. Since $\sigma$ is a bijection from $[n]$ to $[n]$, we see that $A = \mathfrak{A}(\sigma^{-1}(M_{t})) = \mathfrak{A}(M_{t+1})$.
        
        The second possibility is when $|M_{t+1}| = |M_t| + 1$. Then there exists $i\in S_t$ such that $M_{t+1}=M_t\cup\{i\}$. This can happen in exactly two ways: either the chosen move is the dummy move $(i\;i\cdots i)$ or $(i\;a_2\;a_3\cdots\;a_k)$ where $a_2,\cdots,a_k\in M_t$ are distinct. We define $\mathcal{T}_i(M_t)=\{\text{id}\}\cup\{(i\;a_2\;a_3\cdots\;a_k):~a_2,\cdots,a_k\text{ are distinct elements of } M_t\}$ and for $\sigma\in\mathcal{T}_i(M_t)$, the corresponding weight $w(\sigma)=\left\{\begin{array}{ll}C_tq_t, & \sigma =\text{id}  \\q_t, & \sigma \neq \text{id}\end{array}\right.$. Fix $\pi\in \mathfrak{A}(M_{t+1})$, we have 
        \begin{align*}
            \mathbb{P}(\rho_{t+1}=\pi,M_{t+1}=M_t\cup\{i\}|M_t) =\frac{1}{|\mathfrak{A}(M_{t})|}\sum_{\sigma\in\mathcal{T}_i(M)}w(\sigma)\mathbbm{1}_{\sigma^{-1}\pi\in \mathfrak{A}(M_{t})}.
        \end{align*}
%         It is enough to show that it does not depend on $\pi$. If $\pi^{-1}(i)=i$, only dummy move contributes $C_tq_t$ and non-trivial $k$-cycle moves contribute nothing. If $\pi^{-1}(i)=a\in M_t$, in this case, $a_k=a$ is fixed, then there are $(|M_t|-1)_{k-2}$ possible non-trivial $k$-cycle moves can contribute, where each contribute $q_t$. We have shown that
% \begin{itemize}
%     \item If $\pi^{-1}(i)=i$, then the total contribution is $C_tq_t$, coming only from the dummy move. 
%     \item If $\pi^{-1}(i)\in M_t$, then the total contribution is $C_tq_t$, coming from non-trivial $k$-cycles, 
% \end{itemize}
% which completes the proof. 
It is enough to show that
\[
\sum_{\sigma\in \mathcal T_i(M_t)} w(\sigma)\,
\mathbbm{1}_{\sigma^{-1}\pi\in \mathfrak{A}(M_{t})}
\]
does not depend on \(\pi\in \mathfrak A(M_{t+1})\). If \(\pi^{-1}(i)=i\), then the only such \(\sigma\in\mathcal T_i(M_t)\) is
\(\sigma=\mathrm{id}\), so the total contribution is
\[
w(\mathrm{id})=C_t q_t.
\]
If \(\pi^{-1}(i)=a\in M_t\), then \(\sigma\neq \mathrm{id}\) and $\pi(i) = \sigma(i)$. Thus, $\sigma \in \mathcal T_i(M_t)$ must be of the form
\[
\sigma=(i\, \pi(i)\, a_2\,\dots\,a_{k-1}),
\]
where \(a_2,\dots,a_{k-1}\) are distinct elements of \(M_t\setminus\{a\}\).
Hence there are exactly
\[
(|M_t|-1)_{k-2}=C_t
\]
such permutations, and each has weight \(q_t\). So the total contribution is again $C_t q_t$. This completes the proof.
\end{proof}



\begin{lemma}
    With notations defined above, we have $d_{TV}(\mathcal{L}(Z_{\tau_1}),\mathcal{L}(Y_{\tau_2})) \leq n^{-1+o(1)}.$
\end{lemma}

\begin{proof}
For simpler notation, we let $s =|S_t|$ and $|M_t| = n-s.$ Then we first compare the one-step transition rules of the two coupled chains.
\begin{equation}
\begin{aligned}
    &d_{TV}(P_{n},Q_{n,t}) = sC_tq_t
    = \frac{s}{n(n-1)}\frac{(n-s)k}{(n-s+1)}\prod_{i=2}^{k-1}\left(1-\frac{s-1}{n-i}\right)
    = \frac{sk}{n^2}\left(1 + O(s/n)\right),
\end{aligned}
\end{equation}
since
\begin{equation}
\begin{aligned}
    \prod_{i=2}^{k-1}\left(1-\frac{s-1}{n-i}\right) = \left(1 - \frac{k}{n}\right)^{s-1}\exp(O(n^{-1})) = 1 - O((\log n)^{-2})
\end{aligned}
\end{equation}
for $k = o(n/(\log n)^4)$ and $s = O(\log n)$. For any fixed $a >0$, we have $$\{S_{t^*} \geq a\} \subseteq \cup_{A\subseteq[n], |A| = a} \{A \text{ has not been touched up to time }t^*\}.$$ \textcolor{red}{to be removed:}Recall the notation $n_x = n(n-1)(n-2)\cdots (n-x+1)$ and $(n-a)_k = (n-a)(n-a-1)\cdots (n-a-k+1)$. Then
\begin{equation}
\begin{aligned}
    &\Prob\left(|S_{t^*}| \geq  a\right) \leq \binom{n}{a}\left(\frac{\binom{n-a}{k}}{\binom{n}{k}}\right)^{t^*} = \frac{n_a}{a!}\left(\frac{(n-a)_k}{n_k}\right)^{t^*}
    \leq \left(\frac{en}{a}\right)^a \left(1-\frac{a}{n}\right)^{kt^*}\\
    &\leq \left(\frac{en}{a}\right)^a e^{-akt^*/n} = \exp\left( a\left(1 + \log n - \log a - \frac{t^*k}{n}\right) \right).
\end{aligned}
\end{equation}
Recall that $t^*  = \frac{n\log n}{k} + c\frac{n}{k}\log\log\log n + o(\log \log\log n)$ for $|c| < 1/2$.
For $a = \log n$, we have
\begin{equation}
    \Prob\left(S_{t^*} \geq \log n\right) \leq \exp \left( - \log n \log\log n + c\log n \log\log\log n\right) = n^{-\omega(1)}.
\end{equation}
Next, we need to estimate the stopping time.
\begin{equation}\label{eq: step1}
    \begin{aligned}
        &\Prob\left(\max\{ \tau_1, \tau_2\} - t^* \geq \frac{n\log n}{k} \right)\\ 
        &\leq \Prob\left( \tau_1 - t^* \geq \frac{n\log n}{k}, |S_{t^{*}}| \leq a \right) + \Prob\left( \tau_2 - t^* \geq \frac{n\log n}{k}, |S_{t^{*}}| \leq a \right)+\Prob\left(|S_{t^*}| \geq a\right).\\
    \end{aligned}
\end{equation}
We have that 
\begin{equation}\label{eq: step2}
    \begin{aligned}
        &\Prob\left( \tau_1 - t^* \geq \frac{n\log n}{k}, |S_{t^{*}}| \leq a \right) = \Prob\left(\cup_{i \in S_{t^*}} Y_t \text{ does not mark } i \text{ in the next } n\log n/k \text{ steps} , |S_{t^{*}}|
        \leq a \right) \\
        &\leq a \left(1-q_t(n-a)_{k-1}\right)^{n\log n/k} 
        \leq a\left(1-\left(1-\frac{a}{n}\right)^{k-1} \frac{k}{(n-k+1)}\right)^{n\log n/k}\\
        &= a\left(1-\frac{k}{n} + \frac{k^2a}{n^2} + o\left(\frac{1}{n^2}\right)\right)^{n\log n /k}
        = a\exp\left( \left(- \frac{k}{n} + o(n^{-1})\right)n\log n/k\right)\\ &\leq an^{-1 + o(1)}\leq  n^{-1 +o(1)}.
    \end{aligned}
\end{equation}
Same computation for $\tau_2$ gives
\begin{equation}\label{eq: step3}
    \begin{aligned}
        &\Prob\left(Z_t \text{ does not mark }i \text{ in the next } n\log n/k \text{ steps}\right) \leq a\left(1-\frac{(n-a)_{k-1}}{\frac{n!}{(n-k)!k}}\right)^{n\log n/k}\leq n^{-1+o(1)}.
    \end{aligned}
\end{equation}
Then by \eqref{eq: step1}, \eqref{eq: step2}, and \eqref{eq: step3}, we get
\begin{equation}
    \begin{aligned}
        \Prob\left(\tau_1 \neq \tau_2\right) \leq \Prob\left(\max \{\tau_1, \tau_2 \} - t^*> n\log n/k\right) + d_{TV}(P_{n,k},Q_{n,k}) \frac{n\log n}{k} \leq n^{-1+o(1)}.
    \end{aligned}
\end{equation}
Since $Y_{\tau_2} \stackrel{d}{=} \mathfrak{A}_{[n]} \text{ or } \mathfrak{A}_{[n]}^c,$ we get $d_{TV}(\mathcal{L}(Z_{\tau_1}), \mathcal{L}(Y_{\tau_2})) \leq n^{-1+o(1)}$.
\end{proof}
\begin{lemma}\label{lemma::dTV of Z_t^* and X_t^*}
       With notation defined above, we have 
       $d_{TV}(\mathcal{L}(Z_{t^*}),\mathcal{L}(X_{t^*})) \leq (\log n)^{-1/2+o(1)}$.
\end{lemma}
\begin{proof}
    We first compare the joint law of two random pairs $(S_{t^*},X_{t^*})$ and $(S_{t^*}, Z_{t^*})$, 
    \begin{align*}
        d_{TV}(\mathcal{L}(S_{t^*},X_{t^*}),\mathcal{L}(S_{t^*},Z_{t^*})) &= \sum_{S\subset [n]} \Prob(S_{t^*}= S) d_{TV}(\mathcal{L}\left(X_{t^*}~|~S_{t^*}=S\right),\mathcal{L}(Z_{t^*}~|~S_{t^*}=S)) \\
        & = \sum_{S\subset [n]} \Prob(S_{t^*}= S) d_{TV}(\mathcal{L}(X_{t^*}~|~F_{t^*}(S)),\mathrm{Unif}(\Omega_{t^*}(S))) \\
        &\leq \sum_{S\subset [n], |S|\leq \log n} \Prob(S_{t^*}= S) \left(\log n\right)^{-1/2+o(1)} +\Prob\left(|S_{t^*}|>\log n\right) \\
        &\leq \left(\log n\right)^{-1/2+o(1)}+ n^{-\omega(1)}\\
        &=\left(\log n\right)^{-1/2+o(1)}.
    \end{align*}
    Finally, total variation distance cannot increase under marginalization, therefore,
    \begin{align*}
        d_{TV}(\mathcal{L}(Z_{t^*}),\mathcal{L}(X_{t^*})) \leq  d_{TV}(\mathcal{L}(S_{t^*},X_{t^*}),\mathcal{L}(S_{t^*},Z_{t^*})) \leq \left(\log n\right)^{-1/2+o(1)}.
    \end{align*}
\end{proof}

\begin{lemma}
    With notation defined above, we have 
    $$d_{TV}(\mathcal{L}(Z_{\tau_1}),\mathcal{L}(X_{\tau}))\leq \exp\left(-(\log\log n)^{1/2+o(1)}\right).$$
\end{lemma}
\begin{proof}
    % touch and mark steps don't disagree too much.
    Set $T:=\frac{n\log n}{k}$ \JC{We have already use $T$ as a subset of $[n]$. It is better that we don't introduce this $T$ for another meaning}. We have 
    \begin{equation}
        \Prob(\tau_1\neq \tau) \leq \Prob(\tau\leq t^*)+\Prob(|S_t^*|\geq \log n)+\Prob(\tau-t^*\geq T,|S_t^*|\leq \log n)+\Prob(B),
    \end{equation}
    where $B$ is the event that before time $t^*+T$, some chosen $k$-cycle contains at least two elements of $S_{t^*}$ when $|S_{t^*}|\leq \log n$.
    By \Cref{lemma::Self-reduction lemma}, we have 
    \begin{align*}
        \Prob(\tau\leq t^*) &\leq \Prob(\mathcal{F}^{t^*}(\varnothing)) =\left(1+(\log n)^{-1/2+o(1)}\right)\exp\left(-\gamma_{t^*}\right).
    \end{align*}
    Here, we use the fact that
    \begin{align*}
        \gamma_{t^*}&=\exp\left(-\frac{\left(t^*-\lfloor\frac{n\log n}{k}\rfloor\right)\cdot k}{n}\right) \\&\leq \exp\left(\frac{1}{2}\log\log\log n -4\log\log\log\log n+o(1)\right)\\
        &=\left(\log\log n\right)^{\frac{1}{2}+o(1)}.
    \end{align*}
    Therefore, we obtain
    \begin{equation}
        \Prob(\tau \leq t^*)\leq \left(1+(\log n)^{-1/2+o(1)}\right)\exp\left(-(\log\log n)^{1/2}\right).
    \end{equation}
    We have
    \begin{equation}
        \Prob(|S_{t^*}|\geq \log n) \leq n^{-\omega(1)}
    \end{equation}
    Condition on the event $|S_{t^*}|\leq \log n$, fix a point in $S_{t^*}$, it is included in a uniform k-cycle with probability $k/n$. Therefore, over the next $T=n\log n/k$ steps, the probability that this point is still untouched is at most $(1-\frac{k}{n})^T$. Therefore, we have
    \begin{equation}
        \Prob(\tau-t^* >T,\;|S_{t^*}|\leq \log n) \leq |S_{t^*}|\cdot \left(1-\frac{k}{n}\right)^T\leq \log n \cdot\exp\left(-\frac{kT}{n}\right) =n^{-1+o(1)}
    \end{equation}
     Condition on the event $|S_{t^*}|\leq \log n$, starting at $t^*$, in one step transportation, touching scheme and marking scheme differs only when the uniform random $k$-cycle contains at least two points of $S_{t^*}$. The probability of such event is at most $\binom{k}{2}(|S_{t^*}|/n)^2$. Summing over at most $T=n\log n/k$ steps give 
     \begin{equation}
         \Prob(B)\leq \frac{n\log n}{k}\cdot \binom{k}{2}\frac{(\log n)^2}{n^2}\leq \frac{k\cdot(\log n)^3}{n} \leq \frac{1}{\log n}.
     \end{equation}
     Hence, we have
     \begin{equation}
         \Prob(\tau = \tau_1) \geq 1-\exp\left(-(\log\log n)^{1/2+o(1)}\right).
     \end{equation}
By \Cref{lemma::dTV of Z_t^* and X_t^*}, we have 
\begin{align*}
    d_{TV}(\mathcal{L}(Z_{\tau_1}),\mathcal{L}(X_{\tau})) &\leq d_{TV}(\mathcal{L}(Z_{\tau}),\mathcal{L}(X_{\tau}))+\Prob(\tau_1 \neq \tau) \\
    &\leq d_{TV}(\mathcal{L}(Z_{t^*}),\mathcal{L}(X_{t^*})) +\exp\left(-(\log\log n)^{1/2+o(1)}\right) \\
    &\leq \left(\log n\right)^{-1/2+o(1)}+ n^{-\omega(1)} +\exp\left(-(\log\log n)^{1/2+o(1)}\right).
\end{align*}
\end{proof}

\begin{lemma}
    Under the natural coupling in which $(Y_t)_{t\ge t^*}$ and $(\rho_t)_{t\ge t^*}$ both use the same sampled move generated from $Q_{n,t}$ at each time and $Y_{t^*} = \rho_{t^*}$, one has
\[
\tau_2=\tau_3.
\] 
In particular,
    \[
    d_{TV}(Y_{\tau_2}, \rho_{\tau_3}) \leq \frac{1}{\log n}.
    \]
\end{lemma}

\begin{proof}
    Because we have coupled the two chains using the same sampled move $\sigma_t \sim Q_t$ at each time $t \geq t^*$ and set $Y_{t^*} = \rho_{t^*}$. The marked sets for these two chains are identical at any time and therefore $\tau_2 = \tau_3$.   
    Let $T = \frac{n \log n}{k} + t^*$. We know that
    \[
    \Prob\left(\tau_2  \geq T \right) \leq n^{-1+o(1)},\quad \Prob(|S_{t^*}| \geq \log n) \leq n ^{- \omega(1)}.
    \]
    and let $B_T$ be the event that there exists some $1 \leq t \leq T - t^*$ such that $\sigma_{t+t^*}$ contains at least two elements from $[n] \setminus M_{t+t^*}$. Recall that $|[n] \setminus M_{t+t^*}| \leq |[n] \setminus M_{t^*}| = |S_{t^*}|$.
    \[
    \Prob \left( B_T \bigg| |S_{t^*}| \leq \log n \right) \leq \frac{1}{\log n}.
    \]
Thus, by law of total probability, \SC{coupling and original chain}
\[
\begin{split}
    d_{TV}(Y_{\tau_2}, \rho_{\tau_3}) &\leq  \Prob(Y_{\tau_2} \neq \rho_{\tau_3}) \\
    &\leq \Prob(B_{\tau_2})\\
    &\leq \Prob(\tau_2 \geq T) + \Prob \left( B_T \bigg| |S_{t^*}| \leq \log n \right)+ \Prob(|S_{t^*}| \geq \log n)\\
    &\leq \frac{1}{\log n}+n^{-1+o(1)} + n^{-\omega(1)}.
\end{split}
\]
we conclude that $d_{TV}(Y_{\tau_2}, \rho_{\tau_3}) \leq \frac{1}{\log n}$ \JC{is $(1+o(1))\frac{1}{\log n}$ sufficient?}.
\end{proof}


% \begin{lemma}
% Under the natural coupling in which $(Y_t)_{t\ge t^*}$ and $(\rho_t)_{t\ge t^*}$ both use the same sampled move generated from $Q_{n,k}$ at each time and $Y_{t^*} = \rho_{t^*}$, one has
% \[
% \tau_2=\tau_3.
% \]
% Then for every $t\ge t^*$, there exist $\mathfrak A_n$-valued random variables $L_t,R_t$ such that
% \[
% Y_t=L_t\,\rho_t\,R_t.
% \]
% In particular, $Y_{\tau_2}$ is distributed uniformly over $\mathfrak{A}_n$.
% \end{lemma}

% \begin{proof}
% Couple the two chains by using the same sampled move $\sigma_t\sim Q_{n,k}$ at each time $t\ge t^*$. The evolution of the marked set depends only on the sampled move and the current marked set, not on the current permutation. Indeed, if the sampled move is of the form
% \[
% (i_1\cdots i_1), \qquad \text{or} \qquad (i_1\cdots i_k)\ \text{with } i_1\in S_t,\ i_2,\dots,i_k\in M_t,
% \]
% then exactly one new point is marked, namely $i_1$, so
% \[
% M_{t+1}=M_t\cup\{i_1\}.
% \]
% In all other cases, the marked set is merely relabeled:
% \[
% M_{t+1}=\sigma_t^{-1}(M_t).
% \]
% Hence the two chains produce the same process $(M_t)_{t\ge t^*}$, and therefore
% \[
% \tau_2=\tau_3.
% \]

% For the second statement, we argue by induction. At time $t^*$, since $Y_{t^*}=\rho_{t^*}$, the claim holds trivially with $L_{t^*}=R_{t^*}=\mathrm{id}.$ Assume that for some $t\ge t^*$,
% \[
% Y_t=L_t\rho_t R_t.
% \]
% We show that the same holds at time $t+1$. If $\rho_{t+1}=\sigma_t\rho_t$, then the claim holds with
% \[
% L_{t+1}=\sigma_tL_t\sigma_t^{-1},\qquad R_{t+1}=R_t.
% \]
% If $\rho_{t+1}=\sigma_t^{-1}\rho_t\sigma_t$, then the claim holds with
% \[
% L_{t+1}=\sigma_tL_t\sigma_t,\qquad R_{t+1}=\sigma_t^{-1}R_t.
% \]
% By Lemma \ref{lem:uniform}, $\rho_{\tau_2}$ is uniform over $\mathfrak{A}_n$. Thus, $Y_{\tau_2} = L_{\tau_2} \rho_{\tau_2}R_{\tau_2}$ is also uniform over $\mathfrak{A}_n$.
% \end{proof}


\subsection{Even $k$-cycle}
Let $t^* := \lfloor  \frac{n\log n}{k}\rfloor - \frac{n}{2k}\log\log \log n+\frac{4n}{k}\log\log\log \log n$. Let $S_{t^*}$ denote the set of elements in $[n]$ that are not touched by time $t^*$. We consider the following two marking schemes.

\noindent\textbf{Marking scheme A:} at time $t^*$, mark all the elements in $[n] \setminus S_{t^*}$. Let $S_t$ denote the set of all the unmarked elements at time $t$ and let $M_t = [n] \setminus S_t$. 
We define a time-inhomogenous transition kernal as
\[
Q_{n,t}=\begin{cases}
(ab)(i\cdots i) & \text{with probability } \frac{C_tq_t}{\binom{|M_t|}{2}} \ \text{for } i\in S_{t},a,b\in [n]\setminus S_t,\\[6pt]
(i_1 \cdots i_k) & \text{with probability } q_t \ \text{for } i_1 \in S_{t}, i_2, \cdots, i_k \in [n] \setminus S_{t},\\[6pt]
(i_1 \cdots i_k) & \text{with probability } \left( \frac{n!}{k(n-k)!} \right)^{-1} \ \text{otherwise}.
\end{cases}
\]
where $C_t := (|M_t|-1)_{k-2}$ and $q_t := \frac{|M_t|}{|M_t|+1}\left(\frac{n!}{k(n-k)!}\right)^{-1}$.

We update the marked set according to the following rule.
\begin{itemize}
    \item $(ab)(i\cdots i)$ \quad mark $i$ for $i \in S_{t}$.
    \item $(i_1 \cdots i_k)$ mark $i_1$ for $i_1 \in S_t$ and $i_2,\dots, i_k \in [n]\setminus S_t$,
    \item $\sigma=(i_1 \cdots i_k)$ for more than one $j$s such that $i_j\in S_t$. We update the new marked set as follows: we unmark $i_j$ if $\sigma(i_j)\in S_t$ and mark $i_j$ if $\sigma(i_j)\in M_t$. Formally, we update $M_{t+1}=\{i:\sigma(i)\in M_t\} = \sigma^{-1}M_t$ and $S_{t+1}=\{i:\sigma(i)\in S_t\}$
\end{itemize}

\noindent\textbf{Marking scheme B:} at time $t^*$, mark all the elements in $[n] \setminus S_{t^*}$. Let $S_t$ denote the set of all the unmarked elements at time $t$ and let $M_t = [n] \setminus S_t$. At time $t + 1$, we choose a random k-cycle $\sigma = (i_1 \cdots i_k)$ according to the distribution $P_{n}$ defined as follows
\textcolor{red}{to be removed. Repeated}
    \[
P_{n}=\begin{cases}
(i_1 \cdots i_k) & \text{with probability } \left( \frac{n!}{k(n-k)!} \right)^{-1}\\
\text{ otherwise} & \text{with probability } 0.\\[6pt]
\end{cases}
\]
We update the marked set according to the following rule.
\begin{itemize}
    \item If $\sigma=(i_1\cdots i_k)$ with $i_1\in S_t$ and
$i_2,\ldots,i_k\in M_t$, then we mark $i_1$.
\item In all remaining cases, namely when the cycle contains more than one unmarked element, we update the marked set by pulling it back under $\sigma$:
\[
    M_{t+1}
    =
    \{i\in[n]:\sigma(i)\in M_t\}
    =
    \sigma^{-1}(M_t),
\]
and hence
\[
    S_{t+1}
    =
    \{i\in[n]:\sigma(i)\in S_t\}
    =
    \sigma^{-1}(S_t).
\]
\end{itemize}
\noindent\textbf{Coupled Markov processes.}
Let
\[
    \rho_{t^*}=Y_{t^*}=Z_{t^*}
\]
be a common random variable sampled uniformly from $\Omega_{t^*}(S_{t^*})$ . We define four Markov processes
\[
    (X_t)_{t\geq t^*},\qquad
    (Y_t)_{t\geq t^*},\qquad
    (Z_t)_{t\geq t^*},\qquad
    (\rho_t)_{t\geq t^*}
\]
on a common probability space as follows:
\begin{itemize}
    \item The process $(X_t)_{t\geq t^*}$ is the original Markov chain, evolving according to the transition kernel \(P_n\). Let \(\tau\) denote the first time at which all elements have been touched.
    \item The process $(Z_t)_{t\geq t^*}$ is initialized by $Z_{t^*}$, evolving according to the same transition kernel \(P_n\) as $X_t$. Let \(\tau_1\) denote the first time at which all elements have been marked.
    \item The process $(Y_t)_{t\geq t^*}$ is initialized by $Y_{t^*}$, evolving according to the time-inhomogeneous transition kernel \(Q_{n,t}\). Let \(\tau_2\) denote the first time at which all elements have been marked.
    \item Finally, the process $(\rho_t)_{t\geq t^*}$ is initialized by $\rho_{t^*}$, evolving according to the same \(Q_{n,t}\)-updates as $Y_t$, but with a modified action on the current state. More precisely, at time \(t+1\), sample \(\sigma\sim Q_{n,t}\). If either
\[
    \sigma=(i\cdots i)\quad \text{for some } i\in S_t,
\]
or
\[
    \sigma=(i_1\cdots i_k),\qquad i_1\in S_t,\qquad
    i_2,\ldots,i_k\in [n]\setminus S_t,
\]
then set
\[
    \rho_{t+1}:=\sigma\rho_t.
\]
If the sampled move \(\sigma\) satisfies $|\operatorname{supp}(\sigma)\cap S_t|\ge 2,$ then first update the marked and the unmarked set by 
\[
M_{t+1} = \sigma^{-1}(M_t),\quad S_{t+1} = \sigma^{-1}(S_t).
\]
Conditionally on $M_t$, choose an unordered pair $\{U_{t+1},V_{t+1}\}$ uniformly from $M_{t+1}$, independent of $\rho_t$. Then set 
\begin{align*}
    \rho_{t+1} = (U_{t+1}\,V_{t+1})\sigma^{-1}\rho_t\sigma
\end{align*}
Let \(\tau_3\) denote the first time at which all elements have been marked.
\end{itemize}
\begin{lemma}
For $t \geq t^*$, conditioned on the marked set $M_t$, $\rho_t$ is uniform on $\Omega_{t}(S_t)$. In particular, $\rho_{\tau_3}|\tau_3= t$ is uniform on $\Omega_t$. Consequently, we have
\begin{align*}
    \rho_{\tau_3}\sim \Prob(\tau_3\text{ is even})\, U_{\mathfrak{A}_n}+\Prob(\tau_3\text{ is odd})\, U_{\mathfrak{A}_n^c}.
\end{align*}
\end{lemma}
\begin{proof}
    At time $t^*$, this is true by construction. Assume now the claim holds at time $t = t^*$, at time $t^*$, there are two main possibilities. First suppose that $|M_{t+1}|=|M_t|$. There are two subcases. If the sampled \(k\)-cycle \(\sigma\) satisfies $\operatorname{supp}(\sigma)\subset M_t$, then no new element is marked. By the definition of $\rho_t$, $\rho_{t+1} = \sigma \rho_t$, which is a bijection $\Omega_{t}(S_t)\rightarrow \Omega_{t+1}(S_{t+1})$. Hence, $\rho_{t+1}$ is uniform on $\Omega_{t+1}(S_{t+1})$. The second scenario is when the sampled $k$-cycle $\sigma$ satisfies$
|\operatorname{supp}(\sigma)\cap S|\ge2.$ The marked set is updated by $M_{t+1}=\sigma^{-1}(M),~S_{t+1}=\sigma^{-1}(S).$ Set $\eta_{t+1}:=\sigma^{-1}\rho_t\sigma.$ By the induction hypothesis, $\eta_{t+1}$ is uniform on
$\sigma^{-1}\Omega_t(S)\sigma=\Omega_t(\sigma^{-1}(S))=\Omega_t(S_{t+1}).$
By the definition of the even-case update rule, conditionally
on $M_{t+1}$, we choose an unordered pair$ \{U_{t+1},V_{t+1}\}$
uniformly from $\binom{M_{t+1}}{2}$, independently of $\rho_t$, and define
$\rho_{t+1}=(U_{t+1}\,V_{t+1})\eta_{t+1}=(U_{t+1}\,V_{t+1})\sigma^{-1}\rho_t\sigma.$
For every fixed pair $\{u,v\}\subset M_{t+1}$, left multiplication by $(uv)$ is a bijection
$
\Omega_t(S_{t+1})
\longrightarrow
\Omega_{t+1}(S_{t+1}),
$
because \((uv)\) is supported on \(M_{t+1}\) and is an odd permutation. Hence, conditioned
on the chosen pair, \(\rho_{t+1}\) is uniform on \(\Omega_{t+1}(S_{t+1})\). Averaging over the
uniformly chosen pair preserves this uniform law. 

It remains to consider the second possibility: $|M_{t+1}| = |M_t|+1$, then for some $i \in S_t$, $M_{t+1}=M\cup\{i\}, S_{t+1}=S\setminus\{i\}.$ This can happen in exactly two ways. Either the sampled move is a dummy move, or the sampled move is a genuine one-unmarked \(k\)-cycle. Fix $\pi\in\Omega_{t+1}(S_{t+1})$, we have 
\[
\mathbb P(\rho_{t+1}=\pi,\ M_{t+1}=M\cup\{i\}\mid M_t=M)
=
\frac{1}{|\Omega_t(S)|}
\sum_{\alpha}
w(\alpha)\mathbf 1_{\{\alpha^{-1}\pi\in\Omega_t(S)\}},
\] where \(\alpha\) ranges over all dummy transpositions \((ab)\) labelled by \(i\) and all genuine
one-unmarked cycles \((i\,a_2\cdots a_k)\). It is enough to show that $\sum_{\alpha}
w(\alpha)\mathbf 1_{\{\alpha^{-1}\pi\in\Omega_t(S)\}}$ doesn't depend on $\pi$. If $\pi(i) = i$, then only dummy moves can contribute. Indeed, every genuine cycle
\((i\,a_2\cdots a_k)\) sends \(i\) to a marked element, so \(\alpha^{-1}\pi\) cannot fix \(i\).
On the other hand, every dummy transposition \((ab)\), with \(\{a,b\}\subset M\), fixes \(i\)
and gives
$
(ab)^{-1}\pi\in\Omega_t(S).
$
Thus the total contribution is
\[
\binom m2\cdot \frac{C_tq_t}{\binom m2}=C_tq_t.
\]
If \(\pi(i)=a\in M_t\), then no dummy move can contribute, since every dummy transportation fixes $i$. A genuine one-unmarked cycle must satisfy $ \alpha(i)=a.$
Thus \(\alpha\) must be of the form
$\alpha=(i\,a\,a_3\cdots a_k),$
where $a_3,\ldots,a_k$ are distinct elements of $M\setminus\{a\}$. There are exactly $(m-1)_{k-2}=C_t$
such cycles, and each has weight \(q_t\). Hence the total contribution is again $C_tq_t$. Therefore the total preimage weight is independent of $\pi$. It follows that, conditioned on marking \(i\), the random permutation \(\rho_{t+1}\) is uniform on $\Omega_{t+1}(S_{t+1})$. Combining the two cases \(|M_{t+1}|=|M_t|\) and \(|M_{t+1}|=|M_t|+1\), we complete the induction. 
\end{proof}
